\documentclass[10pt,letterpaper]{article}

\usepackage[letterpaper,top=0.45in,bottom=0.45in,left=0.53in,right=0.53in]{geometry}
\usepackage[T1]{fontenc}
\usepackage[utf8]{inputenc}
\usepackage[default,type1]{sourcesanspro}
\usepackage[dvipsnames]{xcolor}
\usepackage{enumitem}
\usepackage{titlesec}
\usepackage{tabularx}
\usepackage{hyperref}
\usepackage{etoolbox}
\input{glyphtounicode}

\pdfgentounicode=1
\providecommand{\ResumeTrack}{general}
\definecolor{primary}{RGB}{0,79,144}
\hypersetup{colorlinks=true,urlcolor=primary,pdftitle={Oswin Franco Cervantes Resume},pdfauthor={Oswin Franco Cervantes}}
\pagestyle{empty}
\setlength{\parindent}{0pt}
\setlength{\tabcolsep}{0pt}
\setlist[itemize]{leftmargin=14pt,topsep=1pt,itemsep=1pt,parsep=0pt,partopsep=0pt}
\titleformat{\section}{\large\bfseries\color{primary}}{}{0pt}{}[\vspace{-3pt}\titlerule]
\titlespacing*{\section}{0pt}{5pt}{3pt}

\newcommand{\role}[4]{%
  \begin{tabularx}{\textwidth}{@{}Xr@{}}
    \textbf{#1} & \textbf{#2} \\
    \textit{#3} & \textit{#4}
  \end{tabularx}\vspace{-4pt}
}

\begin{document}
\fontsize{10pt}{11.75pt}\selectfont

\begin{center}
  {\LARGE\bfseries\color{primary} Oswin Franco Cervantes}\\[-1pt]
  \href{mailto:cervanteso@acm.org}{cervanteso@acm.org} \enspace $|$ \enspace
  \href{tel:+16786222917}{678-622-2917} \enspace $|$ \enspace
  \href{https://www.linkedin.com/in/ocerv/}{linkedin.com/in/ocerv}  \enspace $|$ \enspace
  \href{https://oswincervantes.com/}{oswincervantes.com}
\end{center}
\vspace{-8pt}

\section{Statement of Purpose}
\ifdefstring{\ResumeTrack}{rfic}{%
My goal is to use first-principles physics, electromagnetic modeling, circuit and system-level design, measurement, and computational methods to design and characterize RF and microwave integrated hardware.%
}{\ifdefstring{\ResumeTrack}{photonics}{%
My goal is to use first-principles physics, electromagnetic modeling, circuit and system-level design, microfabrication, measurement, and computational methods to design and characterize photonic integrated hardware.%
}{\ifdefstring{\ResumeTrack}{quantum}{%
My goal is to use first-principles physics, electromagnetic modeling, microwave engineering, circuit and system-level design, microfabrication, measurement, and computational methods to design and characterize quantum hardware.%
}{%
My goal is to use first-principles physics, electromagnetic modeling, circuit and system-level design, measurement, and computational methods to develop integrated hardware spanning semiconductor circuits, RF, photonics, and quantum systems.%
}}}

\section{Education}
\begin{tabularx}{\textwidth}{@{}Xr@{}}
  \textbf{Massachusetts Institute of Technology}, Cambridge, MA & \textbf{Expected May 2028} \\
  B.S. Candidate in Physics and Electrical Engineering with Computing & Sept. 2024--Present
\end{tabularx}
\textit{Selected coursework:} \textbf{Graduate:} High-Frequency Integrated Circuits, Electromagnetics; \textbf{Coursera:} Microwave Engineering and \mbox{Antennas}, RF and Millimeter-Wave Circuit Design; \textbf{EE/Physics:} Solid-State Circuits, Digital Systems Laboratory (Verilog), Semiconductor Electronic Circuits, Signal Processing, Quantum Physics I, Statistical Mechanics; \textbf{Computer Science:} Mathematics for Computer Science, Computational Structures, Low-Level \mbox{Programming}

\vspace{2pt}
\begin{tabularx}{\textwidth}{@{}Xr@{}}
  \textbf{Georgia Tech}, Atlanta, GA --- Dual Enrollment, Mathematics/Computer Science; GPA: 4.0/4.0 & \textbf{Aug. 2022--May 2024}
\end{tabularx}

\section{Integrated Circuit Design Project}
\role{22 nm FinFET Amplifier and Control Chip}{Spring 2026}{MIT Semiconductor Electronic Circuits}{Cambridge, MA}
\begin{itemize}
  \item Designed and simulated a 22 nm FinFET CMOS chip combining an analog amplifier, activation switches, and digital control logic; applied semiconductor device physics and large- and small-signal models.
  \item Created the chip layout in Cadence for a course tapeout and ran Calibre DRC/LVS checks to verify fabrication rules and consistency with the schematic.
  \item Verified through simulation that the amplifier exceeded targets of 26 dB gain at 1 MHz and 300 mV peak-to-peak output swing while operating 48\% below its 500~$\mu$W power limit.
\end{itemize}

\section{Experience}
\role{Undergraduate Researcher --- Prof. William D. Oliver}{Sept. 2025--Present}{MIT Research Laboratory of Electronics, Engineering Quantum Systems (EQuS) Group}{Cambridge, MA}
\begin{itemize}
  \item Own the design and testing of niobium coplanar-waveguide filters for hybrid semiconductor quantum-dot and superconducting devices, minimizing crosstalk, passband insertion loss, and out-of-band noise from room temperature through cryogenic operation.
  \item Create RF layouts in KLayout; model electromagnetic fields in Ansys HFSS/Maxwell using eigenmode and S-parameter analyses plus waveguide, transmission-line, and impedance-matching theory; automate with PyAEDT and extract parasitics with Q3D.
  \item Collaboratively develop CDES, a Python and gdstk framework for superconducting circuit design; contribute reusable components and fabrication GDS layouts through GitHub version control; fabricate niobium devices on silicon by cleanroom photolithography.
\end{itemize}

\role{Avionics Engineer}{Jan. 2026--Present}{MIT Rocket Team}{Cambridge, MA}
\begin{itemize}
  \item Design RF circuits and PCB interconnects in Altium; analyze interconnect behavior with HFSS Driven Terminal and Ansys SIwave; compare simulated S-parameters with vector network analyzer measurements.
\end{itemize}

\role{Hardware Development Engineer (Systems Engineering) Intern}{May 2025--Aug. 2025}{Amazon Lab126, Advanced Products}{Sunnyvale, CA}
\begin{itemize}
  \item Translated ambiguous product requirements with a cross-functional team into system architecture, functional block diagrams, schematics, and PCB layout for a multi-IC solar-charging module compatible with Amazon Kindle devices.
  \item Designed and simulated buck, boost, and buck-boost power circuits in OrCAD Capture and SPICE and laid out PCBs in Cadence Allegro; measured efficiency with oscilloscopes and electronic loads and investigated MPPT.
  \item Implemented and debugged I\textsuperscript{2}C for IC integration; assembled and tested 0402, 0603, and 0804 SMD prototypes and debug boards to validate end-to-end functionality and performance; trained new hires on PCB tools.
\end{itemize}

\section{Technical Skills}
\ifdefstring{\ResumeTrack}{rfic}{%
\textbf{Design:} RF/microwave circuits, solid-state circuits, analog/mixed-signal ICs, digital/RTL design, S-parameters, transmission lines, impedance matching, filters/resonators, SPICE, PCB design\\
\textbf{Layout and Test:} 22 nm FinFET CMOS, custom IC layout, DRC/LVS, SMD assembly, I\textsuperscript{2}C, VNA and oscilloscope measurements\\
\textbf{Software/HDL:} Python, C, Java, Julia, MATLAB, Bash, Verilog, VHDL, Bluespec, RISC-V asm.; NumPy, pandas, gdstk, Git\\
\textbf{EDA Tools:} Ansys (HFSS, Maxwell, Q3D, SIwave, PyAEDT); Cadence (custom IC, AWR Microwave Office, Allegro); Calibre, Altium, OrCAD, KLayout%
}{\ifdefstring{\ResumeTrack}{photonics}{%
\textbf{Design and Analysis:} Electromagnetic modeling, eigenmode analysis, wave propagation, solid-state circuits, digital/RTL design, S-parameters, resonators, signal integrity, analog/RF circuits\\
\textbf{Layout, Fabrication, and Test:} GDS/KLayout, niobium-on-silicon photolithography, cleanroom processing, custom IC layout, DRC/LVS, VNA and oscilloscope measurements\\
\textbf{Software/HDL:} Python, C, Java, MATLAB, Verilog, VHDL; NumPy, pandas, gdstk, Git\\
\textbf{EDA Tools:} Ansys (HFSS, Maxwell, Q3D, SIwave, PyAEDT); Cadence (custom IC, AWR Microwave Office); Calibre, Altium, OrCAD, KLayout%
}{\ifdefstring{\ResumeTrack}{quantum}{%
\textbf{Design and Analysis:} Superconducting circuits, semiconductor quantum dots, solid-state circuits, digital/RTL design, microwave filters/resonators, S-parameters, transmission lines, impedance matching, analog/RF circuits\\
\textbf{Layout, Fabrication, and Test:} GDS/KLayout, niobium-on-silicon photolithography, cleanroom processing, custom IC layout, DRC/LVS, SMD, VNA measurements\\
\textbf{Software/HDL:} Python, C, Java, Julia, Bash, Verilog, VHDL, Bluespec, RISC-V asm.; NumPy, pandas, gdstk, Git\\
\textbf{EDA Tools:} Ansys (HFSS, Maxwell, Q3D, SIwave, PyAEDT); Cadence (custom IC, AWR Microwave Office, Allegro); Calibre, Altium, OrCAD, KLayout%
}{%
\textbf{Design and Analysis:} Electromagnetic modeling, RF/microwave circuits, solid-state circuits, analog/mixed-signal ICs, digital/RTL design, power circuits, S-parameters, transmission lines, impedance matching, SPICE, PCB/system design\\
\textbf{Fabrication and Test:} 22 nm FinFET custom layout, DRC/LVS, GDS/KLayout, niobium-on-silicon photolithography, cleanroom processing, SMD, VNA and oscilloscope measurements\\
\textbf{Software/HDL:} Python, C, Java, Julia, MATLAB, Bash, Verilog, VHDL, Bluespec, RISC-V asm.; NumPy, pandas, gdstk, Git\\
\textbf{EDA Tools:} Ansys (HFSS, Maxwell, Q3D, SIwave, PyAEDT); Cadence (custom IC, AWR Microwave Office, Allegro); Calibre, Altium, OrCAD, KLayout%
}}}

\section{Honors and Certification}
\textbf{Honors:} QuestBridge Scholarship; Amazon Scholar; Jack Kent Cooke Foundation Semifinalist; Stanford Intrologic Scholar\\
\textbf{Certification:} TestOut Network Pro Certification

\end{document}
